% !TEX root = document.tex
\chapter{\label{chap:conclusion}Conclusion}

This thesis presents an algorithm for traversal fusion in strategic programming languages, answering the research question: `Can generic traversal fusion be performed in strategic programming languages, and what is the impact of this optimization on strategic programs?'

To answer the research question, an algorithm that can perform traversal fusion in a strategic programming language was developed based on previous work (TreeFuser and Grafter \cite{SakkaS017, SakkaSN019}). The method relies fundamentally on a dependency analysis of the input program. This dependency analysis can identify opportunities to fuse traversals without altering program results. A proof-of-concept implementation of this algorithm for the Stratego language was implemented and evaluated on several synthetic benchmarks.

We find that traversal fusion can indeed be performed in strategic programming languages but that performance impact in real-world use cases is absent. In some cases, such as when performing full fusion, or when the portion of total runtime taken up by the traversal calls is relatively large, traversal fusion can provide a measurable mean runtime decrease ($>50\%$). However, these situations are unlikely to occur in real programs. In a render tree benchmark simulating a real-world use case, node visits were reduced by over 40\%, but no speedup was measured\footnote{At least, no additional speedup in addition to the unexplained speedup due to code synthesis}.

Compared to traversal fusion for imperative languages, in which traversal fusion provides a significant performance benefit, these results are surprising. The most likely theories to explain the lack of speedup are that fusion introduces some unexpected overhead in Stratego, or that other parts of the program (unrelated to the traversal) take an unexpectedly large amount of time. In either case, we do believe that it is worth investigating these results in more detail, and that it is too early to conclude that the traversal fusion optimization is entirely unviable in strategic programming.

Still, we have shown that traversal fusion can at least be performed on generic traversals in strategic programming languages, laying the foundation for future research in this area.

\section{Future Work}
We make the following suggestions for future work.

First, we recommend that the unexpected lack of performance be further investigated. There are a few different possibilities:
\begin{itemize}
    \item Perform a more thorough investigation into Stratego, the specific code generated by this algorithm, and analyze it in more detail using a profiler
    \item Create more synthetic benchmarks to analyze the performance of different programs, to create a better understanding of which aspects create performance issues
    \item Implement this algorithm in a different strategic programming language to assess whether the unexpected results are specific to Stratego
\end{itemize}

Should it turn out that the performance issues can be resolved and that traversal fusion is indeed a viable optimization for strategic programming, there are several aspects of the algorithm itself that can also be improved. These are, in no particular order:
\begin{itemize}
    \item If the optimization is viable, it would be a good idea to implement a version of the algorithm in the Stratego compiler, that also operates on the full language (not a subset). This would require adding support for additional syntax features.
    \item Currently, rewrite rules are analyzed as a single large statement that cannot be split. Automatically splitting rewrite rules into multiple statements, possibly based on a heuristic, may increase the number of opportunities for fusion.
    \item It may be a good idea to implement optimizations for the algorithm itself. The POC was not developed with performance in mind, and if it is to run as part of the compiler, it would benefit from increased performance.
    \item Assess whether there are better code generation strategies. It is possible that the current matching approach to congruence does not create the fastest possible code. 
    \item Assess whether certain topological orders are better than others. The code generation currently picks a random topological order after fusion. It may turn out that certain traversal orders are simply better.
    \item Attempt to develop heuristics to determine which fusion opportunities are best. The algorithm currently creates every possible fusion, and then picks the fusion with the most amount of fused nodes. This is not a great strategy, as it scales poorly. A good heuristic to pick certain fusion opportunities over others is important to make this optimization practically useful.
    \item Attempt to use results from this algorithm to find new fusion rules that \textit{always} work. We are not certain of how likely this is, but it could be the case that there are certain fusions that are always legal that can be identified using this algorithm. It would at least be interesting to explore this.
\end{itemize}